% Contribution to GT_TOWER_CHECK. Requires amsmath, amssymb, amsthm.
% AI-produced by GPT-6 Astra (OpenAI, ChatGPT Pro mode). No priority claimed.
\subsection{The labelled translation subgroup in the three-division comparison}

The following argument repairs one point in EXIT\_TOWER E.49.  That
proof allows a determinant-minus-one relabelling of the nine translation
labels, and later treats each of the four pair labels as fixed.  The
relabelled action would permute those four pairs.  In the multiplier-one
case under consideration the second pure conjugacy class is excluded by
the peripheral orientation character, as we now prove.  All completions
in this subsection are full profinite completions.

\paragraph{The actual pullback groups.}
Let $T$ be an oriented torus, let
$H=(\mathbb Z/3)^2$ act by translations, and mark its regular orbit
$P=\{p_h:h\in H\}$, with $p_0$ distinguished.  Put
\[
 M=\Gamma(T,P),\qquad
 \widetilde M=\Gamma(T,P\sqcup\{*\})_{\{*\}},
 \qquad \pi:\widetilde M\longrightarrow M.
\]
Here $M$ permits permutations of the nine marks and the subscript on
$\widetilde M$ requires the new mark to be fixed.  The ordinary Birman
sequence has kernel
$\Pi=\pi_1(T\setminus P,*)\cong F_{10}$.
The inverse image $E_H=\pi^{-1}(H)$ is the orbifold fundamental group
of the quotient of $T\setminus P$ by $H$.  Since the action is free,
this is the ordinary fundamental group of the once-punctured quotient
torus.  Consequently the covering gives an exact sequence
\[
 1\longrightarrow F_{10}\longrightarrow E_H\cong F_2
       \xrightarrow{q_H}H\longrightarrow1.                 \tag{CR.1}
\]
The homomorphism $q_H$ is abelianization modulo three followed by an
isomorphism $(\mathbb Z/3)^2\to H$.  Indeed $H$ is abelian of exponent
three and order nine, and $q_H$ is onto.  The profinite Birman sequence
is exact with the full topology on its kernel.  Since $H$ is finite,
the completion of (CR.1) is therefore the actual inverse image of $H$
in the completed universal-curve extension.

Let $F\in\widehat{\mathrm{GT}}$ have multiplier $\lambda=1$.
We use its already constructed action on the ordinary labelled
universal-curve diagram, including the peripheral-power rule.  This is
the tower action used before E.49; no conclusion about the CM
character is used here.  Choose representatives on $M$ and
$\widetilde M$ so that their square with $\pi$ commutes.  To do this,
lift the target inner comparison through the surjective map
$\widehat\pi$ and adjust the representative on
$\widehat{\widetilde M}$ by that lift.  Hence the action restricts to
an isomorphism from $\widehat E_H$ to the inverse image of $F(H)$.

Boggi--Zalesskii Theorem A(i), at $p=3$, realizes $F(H)$ by a discrete
finite subgroup $H'$.  Its conjugator can be chosen pure while allowing
the discrete subgroup $H'$ to vary: subtract a discrete mapping class
with the same nine-label permutation as the original conjugator.
Then $H'$ has exactly the same elementwise label permutations as $H$.
The subgroup $H'$ acts freely on the torus.  In fact a fixed point of
a nonidentity stabilizer is unmarked, since the nine marks form a free
orbit; marking that point produces a nontrivial order-three subgroup
in its universal-curve inverse image.  The original inverse image
$\widehat E_H\cong\widehat F_2$ is torsion-free, and its image under
the tower isomorphism has the same property.  This excludes such a
stabilizer.  Finite realization identifies the resulting free torus
action with a translation action.

Lift the pure conjugator through $\widehat\pi$.  Conjugation by this
lift gives the actual continuous isomorphism
\[
 \theta:\widehat E_H\xrightarrow{\sim}\widehat E_{H'}.
                                                               \tag{CR.2}
\]
Changing the lift changes (CR.2) by an inner automorphism of the
target.  Let $q_{H'}$ identify the deck group with $H$ through its
faithful, unchanged action on the nine labels.  The square on deck
groups then commutes with the identity of $H$:
\[
                       q_{H'}\theta=q_H.                       \tag{CR.3}
\]
This equality follows from preservation of the finite label quotient
and the purity of the correcting conjugator.

\paragraph{The peripheral determinant.}
Write $d_H$ and $d_{H'}$ for positive peripheral loops about the unique
puncture of the two quotient tori.  Every one of the nine upstairs
punctures is a degree-one lift of that puncture.  Its positive loop,
included from the surface kernel into $E_H$, is therefore
$E_H$-conjugate to $d_H$, with exponent one.  The same assertion holds
for $H'$.

The tower action sends each named surface peripheral loop to a
conjugate of its $\lambda$th power in the surface kernel.  Here the
peripheral is the point-pushing loop represented in the universal
curve by the positive twist surrounding the named old mark and the
moving mark.  The labelled universal-curve action preserves that
inertia and its exponent.  Inner comparisons in the base introduce
only the usual surface-peripheral conjugacy: a pure base mapping
class preserves each named peripheral conjugacy class in the fibre.
Thus this is a statement in the fibre group, with its inner ambiguity,
and not merely a conjugacy assertion in the larger mapping class
group.  Applying it to the mark $p_0$ and using the degree-one
inclusion gives
\[
                  \theta(d_H)\sim_{\widehat E_{H'}}
                                     d_{H'}^{\lambda}.           \tag{CR.4}
\]

Choose positive oriented bases $u,v$ and $u',v'$ of the quotient torus
groups.  The same boundary convention identifies the two peripherals
with $[u,v]$ and $[u',v']$, or with their inverses simultaneously.
Let $A\in\operatorname{GL}_2(\widehat{\mathbb Z})$ be the matrix of
$\theta$ on abelianization.  In the quotient by the closed third
lower-central subgroup, expansion of commutators gives
\[
 [u'^{a}v'^{c},u'^{b}v'^{d}]
                   =[u',v']^{ad-bc},\qquad
 \theta([u,v])=[u',v']^{\det A}.                                \tag{CR.5}
\]
The commutator subgroup of this quotient is central, so both
conjugation and the commutator factors in the values of $\theta(u)$
and $\theta(v)$ disappear from this calculation.  The element
$[u',v']$ has infinite profinite order: in
$\operatorname{UT}_3(\mathbb Z/n)$ its image has order $n$, for every
positive integer $n$.  Equations (CR.4)--(CR.5) consequently imply
\[
                             \det A=\lambda=1.                  \tag{CR.6}
\]
This argument uses only the actual peripheral isomorphism (CR.2),
not the CM comparison later in E.49.

Reduce (CR.3) on abelianization modulo three.  The two deck monodromy
isomorphisms satisfy
\[
                  \overline q_{H'}\,(A\bmod3)
                                      =\overline q_H.
                                                               \tag{CR.7}
\]
By (CR.6), their determinants agree.  Oriented quotient-torus basis
changes form $\operatorname{SL}_2(\mathbb Z)$, whose reduction is
onto $\operatorname{SL}_2(\mathbb F_3)$.  The monodromy isomorphisms
form $\operatorname{GL}_2(\mathbb F_3)$, so their two orbits are
exactly their two determinant values.  Equality in (CR.7) therefore
places $H'$ in the same pure conjugacy class as $H$.  The quotient
homeomorphism lifts uniquely after fixing label zero, and that lift
fixes every other label because it intertwines the named deck group.
Thus a pure discrete correcting mapping class suffices.  Together
with the preceding pure profinite conjugator it makes $F(H)=H$
pointwise and preserves all nine labels.

For comparison, the determinant-minus-one coset in
$\operatorname{GL}_2(\mathbb F_3)$ acts by odd permutations on
$\mathbb P^1(\mathbb F_3)$, the four pair labels.  This finite assertion
is checked for all $48$ matrices by the accompanying program.  It
explains why allowing that correction without the determinant
argument does not justify the later assertion that each of the four
forgetting kernels is fixed.

\paragraph{The normalizer topology used in the construction.}
For the translation group, put $N=N_M(H)$ and let $N_0$ fix label zero.
The quotient cover gives
\[
 1\longrightarrow H\longrightarrow N\longrightarrow\Gamma_{1,1}
       \longrightarrow1,
 \qquad N_0\xrightarrow{\sim}\Gamma_{1,1}.
                                                               \tag{CR.8}
\]
Every quotient mapping class preserves the characteristic covering
kernel.  Its lifts form one $H$-torsor, and a unique lift fixes label
zero.  On linear torus representatives the quotient and the map
forgetting the other eight marks induce the same integral homology
matrix.  Thus the second isomorphism in (CR.8) is the actual ordinary
forgetting map.  Every finite quotient of $N_0$ pulls back through
that ambient map.  Hence $N_0$ has the full induced topology.
Its finite-index inclusion in $N$ is induced-open because it
stabilizes a finite label, so $N$ has full induced topology as well.
Boggi--Zalesskii Theorem A(iii), now at $p=3$, gives
\[
                        N_{\widehat M}(H)=\widehat N.           \tag{CR.9}
\]
The genus-one mixed mapping class group is good and of finite virtual
cohomological type: start with the virtually free group
$\Gamma_{1,1}=\operatorname{SL}_2(\mathbb Z)$ and use the full
profinite Birman sequences with finitely generated free kernels.
Filtered unions of finite modules give the arbitrary discrete
continuous module condition in Theorem A, as in the genus-zero
centralizer argument.

\paragraph{Successive fixed-point centralizers.}
The center of $N_0$ contains its unique nonidentity elliptic
involution $\iota$.  The nine original marks consist of one fixed
mark and four pairs.  On adjoining the other fixed points, the
successive orbifold inverse-image groups are
\[
 (C_2*C_2*C_2)*F_4,\quad (C_2*C_2)*F_5,\quad C_2*F_6,\quad F_7.
                                                               \tag{CR.10}
\]
Their even kernels have ranks $10,11,12,13$.  The free-product torsion
calculation, followed by Theorem A at $p=2$, gives respectively
three, two, one, and zero remaining fibre-conjugacy classes of
order-two lifts.  A chosen lift $x$ has centralizer $\langle x\rangle$
in the orbifold inverse image; its intersection with the even kernel
is trivial.

We record the projection comparison explicitly.  In an extension
$1\to K\to E\to L\to1$, the projection of $C_E(x)$ is the stabilizer
of the $K$-conjugacy class of $x$.  An element of that stabilizer has
a lift $e$ with $exe^{-1}=kxk^{-1}$ for some $k\in K$, and
$k^{-1}e$ centralizes $x$.  The kernel is $C_K(x)$, which is trivial
in (CR.10).  Thus the centralizer maps isomorphically to the relevant
finite-index class stabilizer.  Every finite quotient pulls back
through ordinary forgetting to the preimage of that stabilizer in
the ambient group.  This proves its full induced topology and
identifies its completion with the actual completed pullback.
Taking ordered choices and then their finite permutation descent
therefore adjoins the three fixed points functorially.

\paragraph{The final involution quotient.}
After adjoining these points, let $G$ preserve the four fixed labels
individually and preserve the four named pairs, allowing each pair to
flip.  For $C=C_G(\iota)$ the actual unramified double cover gives
\[
 P_0=F_{13}\mathrel{\triangleleft}E_0=F_7,\quad E_0/P_0=C_2,
 \qquad
 1\longrightarrow\langle\iota\rangle\longrightarrow C
       \longrightarrow\Gamma_{0,8}\longrightarrow1.            \tag{CR.11}
\]
At branch cusps the positive upstairs peripheral is the square of
the positive downstairs peripheral; unique roots in the discrete
free group $E_0$ justify descent.  At pair cusps the local degree is
one.  The pure downstairs group preserves the parity homomorphism,
so every such mapping class lifts; the kernel consists of deck
transformations.

Write $P=\widehat P_0$ and $E=\widehat E_0$.  Their extension has
$C_E(P)=1$: a centralizing element has square in $Z(P)=1$, while
$E$ is torsion-free.  Consequently $E$ is the inverse image of the
deck subgroup in $\operatorname{Aut}(P)$.  A centralizing outer
automorphism of $P$ therefore descends to $\operatorname{Out}(E)$;
changing its representative changes the descendant only by an inner
automorphism.  If a class from $\widehat C$ acts trivially on $P$,
adjust its representative by an element of $P$ to make that action
the identity.  Its action on $E$ is then the identity because
$C_E(P)=1$.  Genus-zero congruence faithfulness places the class in
the finite deck kernel of the completion of (CR.11).  That kernel
is detected by the action $-I$ on torus homology modulo three.
Thus $\widehat C$ embeds in the actual ambient $\widehat G$ and has
its full induced topology there.  Direct Theorem A(iii), for the
single order-two subgroup, proves
\[
 C_{\widehat G}(\iota)=\widehat C,\qquad
 \widehat C/\langle\iota\rangle\cong\widehat\Gamma_{0,8}.        \tag{CR.12}
\]
There is no exchange of intersection and closure in this argument.

Each of the four ordinary downstairs marks can be moved once around
a branch mark; its two lifts exchange and the other three pairs stay
fixed.  These four independent monodromies connect all individual
sheet-order choices.  Hence the involution with the specified four
fixed labels and four pairs has one pure conjugacy class.  It can be
straightened by a pure conjugator.

The four maps used by E.49 are now individually named: for
$i=1,2,3,4$, retain the four $E[2]$ branch labels and the $i$th
$E[3]/\{\pm1\}$ pair label and forget the other three pairs.
The determinant repair pins those four pair labels.  The ordered
construction of the three additional order-two lift classes still
permits one common permutation of their names.  We retain this
permutation instead of asserting that an arbitrary correcting
inner automorphism of the original $\Gamma_{1,1}$ fixes their
individual names.

Write the induced eight-point outer action as $(F',\sigma)$ in
the genus-zero decomposition
$\widehat{\mathrm{GT}}\times\mathfrak S_8$.
The four forgetting kernels make $\sigma$ preserve each retained
five-element set, hence fix the four pair labels and preserve the
branch quartet.  The original branch label zero is fixed by the
construction, so the remaining permutation lies in the
$\mathfrak S_3$ of the three newly added branch marks.
On each five-point target E.44 identifies the GT parameter with
$F$; geometric relabelling changes the permutation factor and
does not change that GT factor.  Thus $F'=F$ and the exact
eight-point conclusion is $(F,\sigma)$ with this one common
$\sigma\in\mathfrak S_3$.

The map used subsequently in E.49 retains the original branch
label zero and the four pair labels, and forgets precisely the
other three branch marks.  The preceding $\sigma$ disappears
under this map.  Explicitly, choose a disk containing those three
marks and none of the five retained marks, and represent $\sigma$
by a braid supported in that disk.  After forgetting the three
marks its mapping class is the identity.  Hence the induced
five-point outer action is exactly $(F,1)$, which is the parameter
and labelling required for the monic division-quartic construction.
Equivalently one may remove $\sigma$ at the eight-point stage by
this one geometric relabelling, keeping its induced source inner
comparison.  The proof does not require that all twelve labels
were already pinned by the original translation correction.

Finally this correction is compatible with the single global gauge
in E.267.  After rigidifying the original central involution, write
the corrected source representative as
$\alpha'=\operatorname{Ad}(\gamma)\alpha$ with
$\gamma\in\widehat B$.  For the actual map
$\varphi:\widehat B\xrightarrow{\sim}G_3\subset G_2$,
the inner comparison has the exact algebraic transformation
\[
 \operatorname{Ad}(h)\varphi\alpha'
   =\operatorname{Ad}\bigl(h\varphi(\gamma)\bigr)
       \varphi\alpha.                                      \tag{CR.13}
\]
The multiplier-one source action is the identity on the
abelian quotient $B^{\rm ab}=C_6$, and
$q\varphi(\gamma)$ belongs to the same actual $C_6$ subgroup.
Thus the adjusted comparison still supplies one exponent in
$C_6$.  In particular no independent local signs are introduced
by the remaining branch permutation or by the source inner
ambiguity.
